We first establish two observations that let us analyze MLP fact storage within Transformers. First, MLPs are equivalent to \emph{Hebbian kernel memories} after whitening the empirical feature covariance. Second, MLPs in Transformers need \emph{margins bounded away from zero} (not just positive separability) because attention layers pass noisy queries to MLPs.

\subsection{MLPs Are Hebbian Kernel Memories}
\label{sec:mlps-are-hebbians}

Our first observation is that any MLP can be rewritten as a \emph{Hebbian kernel memory} on its stored examples. The only gap between the plain Hebbian predictor and the original MLP is the empirical feature covariance $\hat{\mSigma}$, so whitening converts any MLP into a (kernel) Hebbian.


\begin{tcolorbox}[breakable,colback=white,colframe=blue!50!black,boxrule=1.5pt,arc=0mm,left=8pt,right=8pt,top=8pt,bottom=8pt,title={\small\bfseries Key Result: MLPs are Hebbian Memories},colbacktitle=blue!50!black,coltitle=white,titlerule=0.5pt]
\begin{theorem}[MLPs as kernel Hebbians, informal]\label{thm:mlp-hebb-whiten}
For stored examples $(\mathbf{x}_{i},\mathbf{y}_{i})$ with
$\mathbf{y}_{i}=\mathrm{MLP}(\mathbf{x}_{i})$ and
$\mathrm{MLP}(\mathbf{x})=\mathbf{B}\phi(\mathbf{x})$
(where for gated MLPs
$\phi(\mathbf{x})=(\mathbf{A}\mathbf{x})\odot \sigma(\mathbf{G}\mathbf{x})$),
define the empirical feature covariance
\[
\hat{\mSigma}:=\frac1F\sum_{i=1}^F
\phi(\mathbf{x}_{i})\phi(\mathbf{x}_{i})^\top .
\]
Assuming $\hat{\mSigma}$ is invertible, define the whitened Hebbian memory
\[
H_{\mathrm{white}}(\mathbf{z})
:=
\frac1F\sum_{i=1}^F
\mathbf{y}_{i}\,K(\mathbf{x}_{i},\mathbf{z})
\]
induced by the kernel
\[
K(\mathbf{x},\mathbf{z})
:=
\phi(\mathbf{x})^\top\hat{\mSigma}^{-1}\phi(\mathbf{z}).
\]
Then
\[
H_{\mathrm{white}}(\mathbf{z})=\mathrm{MLP}(\mathbf{z})
\qquad \text{for all } \mathbf{z}.
\]
Thus, after feature whitening, the MLP is exactly a Hebbian kernel memory.
\end{theorem}
See Appendix~\ref{app:theory:hebbian_mlps:setup} for formal statement and proof. 
\end{tcolorbox}

This reduction motivates the rest of the paper:
constructing an MLP with desirable margin and storage capacity scaling amounts to designing an effective kernel for Hebbian memories.



\subsection{Margins Govern Transformer Usability}
\label{sec:margins-transformer-usability}

Our second observation is that positive margin at stored keys is not enough for a Transformer block to use a fact-storing MLP reliably.
Intuitively, because the attention mechanism perturbs the query before it reaches the MLP, end-to-end usability requires the decoding margin to be bounded away from zero.


\begin{definition}[Synthetic Sequential Factual Recall (SSFR)]
\label{def:ssfr}
Fix a junk-token vocabulary of size $V_{J}$, a model dimension $d$, and a set of $F$
key-value pairs $\{(\mathbf{k}_{i}, \mathbf{v}_{f(i)})\}_{i=1}^{F}$ where
$f : [F] \to [F]$ is a fact set. Then \emph{SSFR} inputs have the form:
$$
  \underbrace{j_{1}, \ldots, j_{J/2}}_{\text{junk prefix}},\;
  \underbrace{k_i}_{\text{key}},\;
  \underbrace{j_{J/2+1}, \ldots, j_{J}}_{\text{junk suffix}},\;
  \underbrace{q}_{\text{query}} \rightarrow \underbrace{v_{f(i)}}_{\text{value}},
$$
where $j_{t} \overset{\mathrm{i.i.d.}}{\sim} \mathrm{Unif}([V_{J}])$ and $q$ is a fixed query token.
The junk length $J$ controls how difficult it is for attention to isolate the relevant key.
\end{definition}

\paragraph{Empirical verification.}
We pretrain an attention-only Transformer block, freeze it, insert a frozen GD-trained fact-storing MLP, and sweep hidden width (see Appendix~\ref{appendix:ssfr-task} for details).
Crucially, although the MLP stores the fact-set as soon as $\gamma_{\min}$ becomes positive, end-to-end SSFR accuracy lags until the margin is bounded away from zero (Figure~\ref{fig:sketched_k2_sweeps}a).
This observation motivates our theoretical study of Hebbian margin bounds in~\Cref{sec:4}.

